\chapter{实数系}
\label{chap:real-number-system}
\label{chap:rational-defects}
\label{chap:ordered-field-completeness}
\label{chap:consequences-completeness}

有理数已经具备四则运算和全序，但某些有界集合在其中没有最小上界，某些彼此越来越接近的有理数列也没有有理极限。为了说明缺少的究竟是什么，需要把域、序和完备性分开考察。本章先规定目标数系应满足的公理，并在“存在一个完备有序域”的条件下推导其结构；\chapterref{chap:dedekind-construction}与\chapterref{chap:cauchy-construction}分别构造这样的域，后者末尾再证明任意两个模型之间存在唯一的保序域同构。因而本章使用符号 \(\mathbb R\) 时，表示这一待构造结构的任一模型，而不是预先假定某一套实数对象已经存在。

\section{有理数系的代数结构及其缺口}

\subsection{从自然数、整数到有理数的代数扩张}


自然数适合计数，但减法不总能在其中进行；整数补入加法逆元，但除法仍可能离开整数。把整数对 \((p,q)\)（\(q\ne0\)）按
\[
 (p,q)\sim(r,s)\quad\Longleftrightarrow\quad ps=rq
\]
分类，并把等价类记为 \(p/q\)，便得到有理数域 \(\Q\)。

\begin{proposition}
有理数集 \(\Q\) 在加、减、乘以及除以非零元素的运算下封闭，并带有与这些运算相容的全序。特别地，若 \(a<b\)，则
\[
 a<\frac{a+b}{2}<b.
\]
\end{proposition}

\begin{proof}
四则运算由
\[
 \frac pq+\frac rs=\frac{ps+rq}{qs},\qquad
 \frac pq\frac rs=\frac{pr}{qs},\qquad
 \left(\frac pq\right)^{-1}=\frac qp\quad(p\ne0)
\]
给出，右端仍由整数对表示；交叉相乘可验证结果不依赖代表元。把分母化为正数后，定义 \(p/q<r/s\) 当且仅当 \(ps<rq\)。整数次序的相容性保证该定义良定，并给出全序。最后，
\[
 \frac{a+b}{2}-a=\frac{b-a}{2}>0,\qquad
 b-\frac{a+b}{2}=\frac{b-a}{2}>0.
\]
\end{proof}

\begin{insightbox}[稠密与完备是两件事]
\(\Q\) 的稠密性只说明任意两个有理数之间还能插入有理数，并未说明每个由有理数逼近确定的位置都已有有理数占据。\(\sqrt2\) 所在的缺口正说明稠密性不能代替完备性。
\end{insightbox}

\subsection{不可公度量与 \texorpdfstring{\(\sqrt2\notin\Q\)}{sqrt(2) 不属于 Q}}


边长为 \(1\) 的正方形对角线若能用同一单位作整数比表示，就会存在正整数 \(p,q\) 满足 \((p/q)^2=2\)。

\begin{theorem}
不存在有理数 \(r\) 满足 \(r^2=2\)。
\end{theorem}

\begin{proof}
反设 \(r=p/q\)，其中 \(p\in\Z\)、\(q\in\Npos\)，且 \(p,q\) 互素。由 \(p^2=2q^2\) 知 \(p^2\) 为偶数，故 \(p\) 为偶数，写成 \(p=2k\)。代回得 \(q^2=2k^2\)，所以 \(q\) 也为偶数，与互素性矛盾。
\end{proof}

实数构造完成以后，\(\sqrt2\) 表示方程 \(x^2=2\) 的正根。上一命题只断言这一方程在 \(\Q\) 中无解，并未预先假定实数系统已经存在。

\begin{remark}
若正整数 \(m\) 不是整数平方，则不存在 \(r\in\Q\) 使 \(r^2=m\)。可比较素因数分解中各素数的指数奇偶性。
\end{remark}

这里遇到的并非四则运算不封闭，而是方程在当前数系中无解。只补入 \(\sqrt2\) 仍会留下 \(\sqrt3\)、三次方程的根以及不由代数方程产生的极限点。因此，逐个添入方程的根不能解决所有由逼近过程产生的缺口。

\begin{proposition}
若素数 \(p\) 整除整数 \(n^2\)，则 \(p\) 整除 \(n\)。因此，若正整数 \(m\) 的素因数分解中有某个指数为奇数，则 \(x^2=m\) 在 \(\Q\) 中无解。
\end{proposition}

\begin{proof}
由算术基本定理，\(n^2\) 的每个素因子指数都是偶数。若 \(p\mid n^2\)，则 \(p\) 在 \(n\) 的分解中指数为正，故 \(p\mid n\)。若 \(r=u/v\) 为最简正有理数且 \(r^2=m\)，则
\[
 u^2=mv^2.
\]
比较任一在 \(m\) 中指数为奇数的素数两边的指数奇偶性，即得矛盾。
\end{proof}

\subsection{有理数轴上的 Dedekind 缺口}


令
\[
 L=\{q\in\Q:q<0\text{ 或 }q^2<2\},\qquad U=\Q\setminus L.
\]
对非负有理数，\(U\) 恰由满足 \(q^2>2\) 的数构成。

\begin{proposition}
集合 \(L,U\) 均非空；每个 \(l\in L\) 都小于每个 \(u\in U\)；并且 \(L\) 没有最大元，\(U\) 没有最小元。
\end{proposition}

\begin{proof}
有 \(-1,1\in L\) 且 \(2\in U\)。若 \(l<0\)，则 \(l<0\le u\)；若 \(l\ge0\)，则 \(l^2<2<u^2\)，故 \(l<u\)。

对 \(q\ge0\) 定义
\[
 T(q)=\frac{2q+2}{q+2}.
\]
直接计算得
\[
 T(q)-q=\frac{2-q^2}{q+2},\qquad
 T(q)^2-2=\frac{2(q^2-2)}{(q+2)^2}.
\]
若 \(q\in L\) 且 \(q\ge0\)，则 \(q<T(q)\) 且 \(T(q)\in L\)；若 \(q<0\)，则 \(q/2\in L\) 且 \(q/2>q\)。故 \(L\) 无最大元。若 \(q\in U\)，则 \(q>0\)、\(q^2>2\)，两式给出 \(T(q)<q\) 且 \(T(q)\in U\)，故 \(U\) 无最小元。
\end{proof}

集合 \(L\) 向下封闭：\(l\in L\) 且 \(r<l\) 蕴含 \(r\in L\)。它与 \(U\) 之间的边界不属于 \(\Q\)。\chapterref{chap:dedekind-construction}将把这类左部本身作为新数。

\subsection{在 \texorpdfstring{\(\Q\)}{Q} 中有界却没有上确界的集合}


\begin{definition}
设 \(A\subseteq\Q\)。若 \(M\in\Q\) 满足 \(a\le M\) 对所有 \(a\in A\) 成立，则 \(M\) 是 \(A\) 在 \(\Q\) 中的上界。若上界 \(s\in\Q\) 还不大于任何其他有理上界，则称 \(s\) 为 \(A\) 在 \(\Q\) 中的上确界。
\end{definition}

\begin{theorem}
集合
\[
 S=\{q\in\Q:q\ge0,\ q^2<2\}
\]
非空且在 \(\Q\) 中有上界，但没有 \(\Q\) 中的上确界。
\end{theorem}

\begin{proof}
\(1\in S\)，而 \(2\) 是上界。反设 \(s\in\Q\) 是上确界。由 \(S\) 含正数知 \(s>0\)，且上一节排除了 \(s^2=2\)。

若 \(s^2<2\)，则 \(T(s)>s\) 且 \(T(s)\in S\)，与 \(s\) 为上界矛盾。若 \(s^2>2\)，则 \(0<T(s)<s\) 且 \(T(s)^2>2\)。若 \(x\in S\) 且 \(x\ge T(s)\)，非负数平方保持次序，从而 \(x^2\ge T(s)^2>2\)，矛盾。因此 \(T(s)\) 是小于 \(s\) 的上界，与 \(s\) 的最小性矛盾。
\end{proof}

“有上界”只要求至少一个控制量；“有上确界”还要求全体上界中存在最小者。两者不能混同。

\begin{proposition}[同一缺口的两侧]
集合 \(S\) 在 \(\Q\) 中的上界全体恰为
\[
 \{u\in\Q:u>0,\ u^2>2\}.
\]
因此“\(S\) 没有上确界”与“右部 \(U\) 没有最小元”是同一事实的两种表述。
\end{proposition}

\begin{proof}
若 \(u\) 是上界，则 \(u\ge1>0\)。若 \(u^2<2\)，变换 \(T(u)\) 给出大于 \(u\) 的 \(S\) 中元素，矛盾；等号又不可能，故 \(u^2>2\)。反之，若 \(u>0\)、\(u^2>2\)，则每个 \(s\in S\) 都满足 \(s^2<2<u^2\)，由非负数平方的严格保序性得 \(s<u\)，所以 \(u\) 是上界。
\end{proof}

\subsection{在 \texorpdfstring{\(\Q\)}{Q} 中没有极限的 Cauchy 列}


\begin{definition}
有理数列 \((a_n)\) 称为有理 Cauchy 列，如果
\[
\forall\varepsilon\in\Q,\ \varepsilon>0,\ \exists N\in\Npos,\
 \forall m,n\ge N,\quad |a_m-a_n|<\varepsilon.
\]
若存在 \(a\in\Q\)，使对每个正有理数 \(\varepsilon\)，从某项起均有 \(|a_n-a|<\varepsilon\)，则称它在 \(\Q\) 中收敛到 \(a\)。
\end{definition}

对 \(n\in\Npos\)，令 \(a_n=k_n/10^n\)，其中 \(k_n\) 是满足 \(k_n\ge0\) 且 \((k_n/10^n)^2<2\) 的最大整数。候选整数至少含 \(0\)，并且均小于 \(2\cdot10^n\)，所以最大值存在。

\begin{theorem}
数列 \((a_n)\) 是有理 Cauchy 列，但不在 \(\Q\) 中收敛。
\end{theorem}

\begin{proof}
由最大性，
\[
 a_n^2<2\le(a_n+10^{-n})^2.
\]
若 \(m\ge n\)，则 \(a_n\) 也能写成分母为 \(10^m\) 的候选数，故 \(a_m\ge a_n\)；而 \(a_m\ge a_n+10^{-n}\) 将推出 \(a_m^2\ge2\)。所以
\[
 0\le a_m-a_n<10^{-n}.
\]
给定正有理数 \(\varepsilon\)，直接按其分数表示选取 \(N\) 使 \(10^{-N}<\varepsilon\)，便得 Cauchy 性。

反设 \(a_n\to a\in\Q\)。由 \(0\le a_n<2\)，三角不等式给出
\[
 |a_n^2-a^2|=|a_n-a||a_n+a|\longrightarrow0.
\]
另一方面，
\[
 0<2-a_n^2\le 2a_n10^{-n}+10^{-2n}
 <4\cdot10^{-n}+10^{-2n},
\]
故 \(a_n^2\) 逼近 \(2\)。于是对任意正有理数 \(\varepsilon\)，都有 \(|a^2-2|<\varepsilon\)。若 \(a^2\ne2\)，取 \(\varepsilon=|a^2-2|/2\) 即矛盾。因此 \(a^2=2\)，与第二节定理矛盾。
\end{proof}

十进制并非关键。以任意整数 \(b\ge2\) 为基数，令
\[
 a_n^{(b)}=\frac{k_n}{b^n},
\]
其中 \(k_n\) 是使 \((k_n/b^n)^2<2\) 的最大非负整数，便有
\[
 0\le a_m^{(b)}-a_n^{(b)}<b^{-n}\qquad(m\ge n).
\]
所有这些列都从缺口左侧逼近同一边界。若取两个基数构造的列 \((a_n^{(b)})\)、\((a_n^{(c)})\)，可由各自网格宽度证明它们的差趋于零；\chapterref{chap:cauchy-construction}将把这种“差为零”的列识别为同一个新数。

\begin{proposition}
有理 Cauchy 列一定有界，但有界有理数列未必是 Cauchy 列；相邻项之差趋于零也不足以推出 Cauchy 性。
\end{proposition}

\begin{proof}
对 Cauchy 列在定义中取 \(\varepsilon=1\)，尾部由某一固定项加 \(1\) 控制，前有限项再取绝对值最大值。数列 \((-1)^n\) 有界但不是 Cauchy 列。调和部分和
\[
 H_n=\sum_{k=1}^n\frac1k
\]
满足 \(H_{n+1}-H_n\to0\)，但例如
\[
 H_{2n}-H_n=\sum_{k=n+1}^{2n}\frac1k\ge\frac12,
\]
故不是 Cauchy 列。
\end{proof}

\section{有序域、绝对值与区间}

\subsection{域的代数结构}


\begin{definition}
集合 \(F\) 连同运算 \(+\) 与 \(\cdot\) 称为一个\term{域}，如果加法、乘法分别满足交换律和结合律，乘法对加法满足分配律，并且存在互异元素 \(0,1\in F\)，使每个 \(x\in F\) 有加法逆元 \(-x\)，每个非零 \(x\in F\) 有乘法逆元 \(x^{-1}\)。
\end{definition}

减法定义为 \(x-y=x+(-y)\)，除法定义为 \(x/y=xy^{-1}\)（\(y\ne0\)）。

\begin{proposition}
在域中，两个单位元和每个逆元都是唯一的；此外
\[
 0x=0,\qquad (-x)y=-(xy),\qquad (-x)(-y)=xy.
\]
若 \(xy=0\)，则 \(x=0\) 或 \(y=0\)。
\end{proposition}

\begin{proof}
若 \(0,0'\) 都是加法单位元，则 \(0=0+0'=0'\)；其他唯一性同理。由
\[
 0x=(0+0)x=0x+0x
\]
消去 \(0x\) 得 \(0x=0\)，其余符号法则由逆元唯一性推出。若 \(xy=0\) 且 \(x\ne0\)，两边乘 \(x^{-1}\) 得 \(y=0\)。
\end{proof}

\(\Q,\R,\C\) 都是域，\(\Z\) 不是域。本章讨论还带有相容全序的域；\(\C\) 不存在这种全序。

\subsection{序与域运算的相容性}


\begin{definition}
域 \(F\) 上的关系 \(\le\) 称为与域运算相容的\term{全序}，如果它满足自反性、反对称性、传递性和任意两元素可比较，并且
\begin{enumerate}
 \item \(x\le y\Rightarrow x+z\le y+z\)；
 \item \(0\le x,\ 0\le y\Rightarrow0\le xy\)。
\end{enumerate}
带有这种全序的域称为\term{有序域}。
\end{definition}

\begin{proposition}
在有序域中，\(1>0\)；对 \(x\ne0\)，有 \(x^2>0\)；若 \(0<x<y\)，则 \(0<y^{-1}<x^{-1}\)。
\end{proposition}

\begin{proof}
若 \(1<0\)，则 \(-1>0\)，从而 \(1=(-1)^2>0\)，矛盾，故 \(1>0\)。对非零 \(x\)，\(x\) 与 \(-x\) 中恰有一个为正，故 \(x^2>0\)。若 \(x>0\)，其逆元也为正；由 \(x<y\) 乘正数 \(x^{-1}y^{-1}\)，得到 \(y^{-1}<x^{-1}\)。
\end{proof}

有序域的特征为零，因为每个有限和 \(n\cdot1\) 都为正。因此其中含有 \(\Q\) 的一个保持运算与次序的副本。

\subsection{有序域中的区间、绝对值与基本不等式}


对 \(a<b\) 定义通常的开、闭、半开区间。绝对值定义为
\[
 |x|=\begin{cases}x,&x\ge0,\\-x,&x<0.\end{cases}
\]

\begin{theorem}[绝对值不等式]
对任意 \(x,y\in F\)，
\[
 |xy|=|x||y|,\qquad |x+y|\le|x|+|y|,\qquad
 \bigl||x|-|y|\bigr|\le|x-y|.
\]
\end{theorem}

\begin{proof}
乘积式按符号分类。由 \(-|x|\le x\le|x|\) 与对应的 \(y\) 不等式相加，得到
\[
 -(|x|+|y|)\le x+y\le|x|+|y|,
\]
即三角不等式。由 \(x=(x-y)+y\) 得 \(|x|-|y|\le|x-y|\)；交换 \(x,y\) 并合并两式，得到反三角不等式。
\end{proof}

对 \(r>0\)，\(|x-a|<r\) 等价于 \(a-r<x<a+r\)。后续极限定义将用这一关系把误差估计改写为区间邻域。

\section{上界、最大元、上确界与确界公理}

\subsection{上界、下界、最大元与最小元}


\begin{definition}
设 \(A\subseteq F\)。若 \(M\in F\) 满足 \(a\le M\) 对所有 \(a\in A\) 成立，则 \(M\) 是 \(A\) 的\term{上界}。若 \(m\in A\) 且 \(a\le m\) 对所有 \(a\in A\) 成立，则 \(m\) 是 \(A\) 的\term{最大元}。下界、最小元对称定义。
\end{definition}

最大元与最小元若存在则唯一。空集的每个域元素都是上界和下界，但空集没有最大元或最小元。确界公理因此显式要求集合非空。

\begin{example}
区间 \((0,1)\) 有上界 \(1\)，却没有最大元：若 \(0<x<1\)，则 \(x<(x+1)/2<1\)。区间 \((0,1]\) 的最大元为 \(1\)。
\end{example}

记上界集为
\[
 \operatorname{UB}(A)=\{M\in F:\forall a\in A,\ a\le M\}.
\]
若 \(\operatorname{UB}(A)\) 非空，它向上封闭：\(M\) 是上界且 \(M'\ge M\) 时，\(M'\) 仍是上界。类似地，下界集 \(\operatorname{LB}(A)\) 向下封闭。上确界若存在，就是上界集的最小元；下确界则是下界集的最大元。这个改写把确界问题转化为“界的集合有没有端点”。

\begin{proposition}
设 \(A,B\subseteq F\) 非空。
\begin{enumerate}
 \item 若 \(A\subseteq B\)，则 \(\operatorname{UB}(B)\subseteq\operatorname{UB}(A)\)；
 \item \(M\) 是 \(A\cup B\) 的上界，当且仅当它同时是 \(A,B\) 的上界；
 \item 若 \(A\) 有最大元，则该最大元是上确界；但上确界不必属于 \(A\)。
\end{enumerate}
\end{proposition}

\begin{proof}
前两项直接展开全称量词。第三项中，最大元属于 \(A\) 且控制全部元素，所以是上界；任意上界还必须不小于这个属于 \(A\) 的元素，故它是最小上界。开区间给出最后一句的反例。
\end{proof}

\subsection{上确界与下确界}


\begin{definition}
设非空集合 \(A\subseteq F\) 有上界。若 \(s\) 是 \(A\) 的上界，并且 \(s\le M\) 对每个上界 \(M\) 成立，则称 \(s\) 为 \(A\) 的\term{上确界}，记为 \(s=\sup A\)。下确界 \(\inf A\) 对称定义。
\end{definition}

\begin{theorem}[上确界的近似刻画]
\(s=\sup A\) 当且仅当
\begin{enumerate}
 \item \(a\le s\) 对所有 \(a\in A\) 成立；
 \item 对每个 \(\varepsilon>0\)，存在 \(a_\varepsilon\in A\) 使 \(s-\varepsilon<a_\varepsilon\le s\)。
\end{enumerate}
\end{theorem}

\begin{proof}
若第二项对某个 \(\varepsilon>0\) 失败，则 \(s-\varepsilon\) 是更小的上界，矛盾。反之，第一项说明 \(s\) 是上界。若有上界 \(M<s\)，取 \(\varepsilon=s-M\)，第二项给出 \(a_\varepsilon>M\)，仍矛盾。
\end{proof}

\begin{proposition}
若相应确界存在，则
\[
 \inf A=-\sup(-A),\qquad
 \sup(A+c)=\sup A+c,\qquad
 \sup(\lambda A)=\lambda\sup A\quad(\lambda>0).
\]
\end{proposition}

\begin{proof}
下界在取相反数后对应上界；平移和正数乘法则分别给出保持大小关系的双射。逐项核对上界性与最小性即可。
\end{proof}

\begin{proposition}[误差刻画的变体]
若 \(s\) 是 \(A\) 的上界，则以下条件等价：
\begin{enumerate}
 \item \(s=\sup A\)；
 \item 对每个 \(t<s\)，存在 \(a\in A\) 使 \(t<a\)；
 \item 对每个正数 \(\eta\)，都有 \(A\cap(s-\eta,s]\ne\varnothing\)。
\end{enumerate}
\end{proposition}

\begin{proof}
第一项与第二项的等价只需在上确界近似刻画中令 \(\varepsilon=s-t\)。第二项说明任何小于 \(s\) 的数都不是上界，也正是“\(s\) 是最小上界”。第三项是同一陈述的邻域形式。
\end{proof}

\subsection{确界公理与完备有序域}


\begin{definition}
有序域 \(F\) 称为\term{完备有序域}，如果它的每个非空且有上界的子集都在 \(F\) 中有上确界。
\end{definition}

由取相反数可知，下确界原理与确界公理等价。\chapterref{chap:real-number-system}的集合 \(S\subseteq\Q\) 说明 \(\Q\) 不完备。

\begin{theorem}[正数平方根]
在完备有序域 \(F\) 中，对每个 \(a>0\)，存在唯一的 \(s>0\) 使 \(s^2=a\)。
\end{theorem}

\begin{proof}
令 \(A=\{x\in F:x\ge0,\ x^2\le a\}\)。它含 \(0\)，且 \(a+1\) 是上界。令 \(s=\sup A\)。因 \(a/(a+1)\in A\)，故 \(s>0\)。

若 \(s^2<a\)，令 \(d=a-s^2>0\)，取
\[
 0<h<\min\left\{1,\frac d{2s+1}\right\}.
\]
则 \((s+h)^2-s^2<(2s+1)h<d\)，所以 \(s+h\in A\)，与上界性矛盾。若 \(s^2>a\)，令 \(d=s^2-a>0\)，取
\[
 0<h<\min\left\{s,\frac d{2s}\right\}.
\]
则 \((s-h)^2>a\)。任何 \(x\in A\) 都不能大于 \(s-h\)，故 \(s-h\) 是更小的上界，矛盾。因此 \(s^2=a\)。

若 \(t>0\) 且 \(t^2=s^2\)，则 \((t-s)(t+s)=0\)。因 \(t+s>0\)，有 \(t=s\)。
\end{proof}

确界公理只用于产生候选 \(s\)；两个有理扰动分别排除 \(s^2<a\) 与 \(s^2>a\)。这填补了\chapterref{chap:real-number-system}由 \(x^2=2\) 产生的缺口。

\subsection{上确界与最大值的区别}


若 \(A\) 有最大元 \(m\)，则 \(m=\sup A\)。逆命题不成立：\((0,1)\) 的上确界是 \(1\)，但 \(1\notin(0,1)\)。

\begin{proposition}
非空有上界集合 \(A\) 有最大元，当且仅当 \(\sup A\in A\)；此时最大元等于上确界。
\end{proposition}

\begin{proof}
最大元既是上界，又不大于任何上界，所以等于上确界。反之，属于 \(A\) 的上确界支配 \(A\) 的全部元素，因而是最大元。
\end{proof}

\begin{example}[边界位置与边界归属]
令
\[
 A=(0,1),\qquad B=(0,1/2),\qquad C=(0,1]\!.
\]
前两个集合都没有最大元，但上确界分别为 \(1\) 与 \(1/2\)；集合 \(A,C\) 有同一个上确界，只有 \(C\) 包含边界并具有最大元。确界记录边界位置，最大元还记录边界是否属于集合。
\end{example}

\begin{proposition}[确界的稳定比较]
设 \(A,B\subseteq F\) 非空且有上界。若对每个 \(a\in A\) 都存在 \(b\in B\) 使 \(a\le b\)，则
\[
 \sup A\le\sup B.
\]
\end{proposition}

\begin{proof}
\(\sup B\) 控制每个 \(b\in B\)，因而按假设也控制每个 \(a\in A\)，所以它是 \(A\) 的上界。最小上界性给出结论。这个命题常用于比较删去有限项、平移或逐点放大后的集合边界。
\end{proof}

下面几条运算公式以后会反复使用。它们的证明形式相似，但符号条件不能省略。

\begin{proposition}[确界的代数运算]
设 \(A,B\subseteq F\) 非空且有界，\(c,\lambda\in F\)。则
\[
 \sup(A+c)=\sup A+c,
\]
并且
\[
 \sup(\lambda A)=
 \begin{cases}
   \lambda\sup A,&\lambda>0,\\
   0,&\lambda=0,\\
   \lambda\inf A,&\lambda<0.
 \end{cases}
\]
此外，若 \(A+B=\{a+b:a\in A,b\in B\}\)，则
\[
  \sup(A+B)=\sup A+\sup B,\qquad
  \inf(A+B)=\inf A+\inf B.
\]
\end{proposition}

\begin{proof}
平移是保序双射，正数伸缩保序，负数伸缩反序，因此前三个公式直接由“上界”与“最小上界”的定义得到。以下证明 Minkowski 和的上确界公式。数 \(s=\sup A+\sup B\) 显然控制每个 \(a+b\)，所以是上界。给定 \(\varepsilon>0\)，分别取
\[
  a>\sup A-\varepsilon/2,\qquad
  b>\sup B-\varepsilon/2.
\]
于是 \(a+b>s-\varepsilon\)，故任何小于 \(s\) 的数都不能是 \(A+B\) 的上界。下确界公式可对 \(-A,-B\) 使用上确界公式。
\end{proof}

若要求积集 \(AB=\{ab:a\in A,b\in B\}\) 的确界，则必须另加符号条件；例如 \(A,B\subseteq[0,\infty)\) 时才有 \(\sup(AB)=\sup A\sup B\)。负数会反转次序，跨过零的集合还要同时比较四个端点。确界公式不能仅凭外形类推。

\section{Archimedes 性质、整数部分与稠密性}

\subsection{Archimedes 性质及其证明}


\begin{theorem}[Archimedes 性质]
每个满足确界公理的有序域 \(F\) 都是 Archimedes 的：对每个 \(x\in F\)，存在 \(n\in\N\) 使 \(n>x\)。
\end{theorem}

\begin{proof}
反设自然数在 \(F\) 的有理副本中有上界，令 \(s=\sup\N\)。数 \(s-1\) 不是上界，所以存在 \(n\in\N\) 使 \(n>s-1\)。于是 \(n+1>s\)，而 \(n+1\in\N\)，与 \(s\) 为上界矛盾。
\end{proof}

证明的实质性步骤是取得 \(\sup\N\)；其余只用自然数对后继封闭。在非 Archimedes 有序域中可能存在大于每个自然数的元素，所以该结论不是有序域公理的形式后果。

\begin{proposition}[Archimedes 性质的等价表述]
在有序域中，以下命题等价：
\begin{enumerate}
 \item 对每个 \(x\)，存在 \(n\in\N\) 使 \(n>x\)；
 \item 对每个 \(\varepsilon>0\)，存在 \(n\in\Npos\) 使 \(1/n<\varepsilon\)；
 \item 对任意 \(x>0,y>0\)，存在 \(n\in\N\) 使 \(nx>y\)。
\end{enumerate}
\end{proposition}

\begin{proof}
第一项用于 \(1/\varepsilon\) 即得第二项。第二项用于 \(\varepsilon=x/y\)，取倒数并整理得第三项。第三项令 \(x=1\)，再把任意非正 \(y\) 单独处理，即得第一项。
\end{proof}

因此 Archimedes 性质排除了非零无穷小和正无穷大元素：若 \(\eta>0\) 且 \(\eta<1/n\) 对所有 \(n\) 成立，第二项会对 \(\varepsilon=\eta\) 产生矛盾；若 \(H>n\) 对所有 \(n\) 成立，则第一项对 \(H\) 失败。

\subsection{自然数无上界与任意小正数的存在}


\begin{corollary}
若 \(\varepsilon>0\)，则存在 \(n\in\Npos\) 使
\[
 0<\frac1n<\varepsilon.
\]
更一般地，对 \(a>0\) 存在 \(n\in\Npos\) 使 \(n\varepsilon>a\)。
\end{corollary}

\begin{proof}
对 \(1/\varepsilon\) 应用 Archimedes 性质，取 \(n>1/\varepsilon\)，再利用正数取倒数反向。第二式对 \(a/\varepsilon\) 应用同一定理。
\end{proof}

“任意小正数”表示对每个预先给定的正阈值都能找到更小的正数，并不存在一个大于零且小于所有正实数的固定实数。

\begin{corollary}[网格可达到任意精度]
给定区间长度 \(\ell>0\) 与目标误差 \(\varepsilon>0\)，存在 \(n\in\Npos\) 使
\[
 \frac{\ell}{n}<\varepsilon.
\]
因此有限区间可以用有限个长度小于 \(\varepsilon\) 的等长小区间覆盖。
\end{corollary}

\begin{proof}
取 \(n>\ell/\varepsilon\)。若以步长 \(\ell/n\) 从左端点划分，共需 \(n\) 个小区间。
\end{proof}

\subsection{整数部分定理}


\begin{theorem}[整数部分定理]
对每个 \(x\in\R\)，存在唯一 \(m\in\Z\) 使
\[
 m\le x<m+1.
\]
该整数记为 \(\lfloor x\rfloor\)。
\end{theorem}

\begin{proof}
由 Archimedes 性质，集合 \(A=\{n\in\N:n>x\}\) 非空。取其最小元 \(n_0\)。若 \(x\ge0\)，最小性给出 \(n_0-1\le x<n_0\)，可取 \(m=n_0-1\)。若 \(x<0\)，对 \(-x\) 取正整数 \(N>-x\)，于是 \(-N<x\)。集合
\[
 B=\{k\in\Z:k\le x\}
\]
非空，并被任意整数 \(n_0>x\) 从上方控制。把有限区间 \([-N,n_0]\cap\Z\) 中属于 \(B\) 的整数取最大者，得到 \(m\)，其最大性给出 \(x<m+1\)。

若 \(m,m'\) 都满足条件且 \(m<m'\)，则整数差给出 \(m+1\le m'\le x\)，与 \(x<m+1\) 矛盾。故唯一。
\end{proof}

等价地，
\[
 \lfloor x\rfloor\le x<\lfloor x\rfloor+1,\qquad
 \{x\}:=x-\lfloor x\rfloor\in[0,1),
\]
其中 \(\{x\}\) 称为小数部分。

\begin{proposition}[整数部分的基本运算]
对 \(x,y\in\R\)、\(m\in\Z\)，有
\[
 \lfloor x+m\rfloor=\lfloor x\rfloor+m,
\]
\[
 \lfloor x\rfloor+\lfloor y\rfloor
 \le\lfloor x+y\rfloor
 \le\lfloor x\rfloor+\lfloor y\rfloor+1.
\]
上式右侧的 \(1\) 不能统一删去。
\end{proposition}

\begin{proof}
第一式把
\(\lfloor x\rfloor\le x<\lfloor x\rfloor+1\)
整体加 \(m\)，再用唯一性。写
\[
 x=\lfloor x\rfloor+\{x\},\qquad
 y=\lfloor y\rfloor+\{y\},
\]
且 \(0\le\{x\}+\{y\}<2\)，便得第二组不等式。取 \(x=y=3/4\) 时右侧修正项确实出现。
\end{proof}

\subsection{有理数在实数中的稠密性}


\begin{theorem}
若 \(x<y\)，则存在 \(q\in\Q\) 使 \(x<q<y\)。
\end{theorem}

\begin{proof}
取 \(n\in\Npos\) 使 \(1/n<y-x\)。令 \(m=\lfloor nx\rfloor+1\)。于是
\[
 nx<m\le nx+1<ny.
\]
除以正整数 \(n\)，得到 \(x<m/n<y\)。
\end{proof}

重复在较小区间中应用该定理可得任意开区间含无穷多个有理数。稠密性不意味着有理数“占据大多数”实数；\chapterref{chap:mathematical-language}已证明 \(\Q\) 可数无限，本章末将证明 \(\R\) 不可数。

\begin{proposition}[定量有理逼近]
对任意 \(x\in\R\) 与 \(n\in\Npos\)，存在 \(m\in\Z\) 使
\[
 \frac mn\le x<\frac{m+1}{n}.
\]
特别地，存在分母为 \(n\) 的有理数 \(q\) 使
\[
 \abs{x-q}<\frac1n.
\]
\end{proposition}

\begin{proof}
取 \(m=\lfloor nx\rfloor\)，第一式由整数部分定理除以 \(n>0\) 得到。可取 \(q=m/n\)；若采用最近网格点，还可把误差改进到 \(1/(2n)\)，但中点处需要约定取哪一侧。
\end{proof}

\subsection{无理数在实数中的稠密性}


\begin{theorem}
若 \(x<y\)，则存在无理数 \(\alpha\) 使 \(x<\alpha<y\)。
\end{theorem}

\begin{proof}
先取 \(q\in\Q\cap(x,y)\)。由任意小正数性质，取 \(n\in\Npos\) 使 \(\sqrt2/n<y-q\)。则
\[
 x<q<q+\frac{\sqrt2}{n}<y.
\]
若 \(q+\sqrt2/n\) 为有理数，减去 \(q\) 再乘 \(n\) 将推出 \(\sqrt2\in\Q\)，矛盾。
\end{proof}

因此有理数与无理数都在每个非空开区间中出现。两者在序上的稠密程度相同，可数性却不同。

\section{十进制表示、不可数性与完备性的需求}

\subsection{分析学对数系提出的完备性要求}


前面已经用三种方式看见同一个缺口：分割没有有理边界，有界集合没有有理上确界，Cauchy 列没有有理极限。\chapterref{chap:equivalent-completeness}将证明，在实数中相应的三条完备性命题彼此等价。

仅补入 \(x^2=2\) 的正根并不充分，还要处理其他代数方程以及一般逼近过程产生的缺口。所需扩张应保留 \(\Q\) 的有序域运算，填补全部 Dedekind 缺口，并使内部 Cauchy 逼近具有极限。本章把这些要求公理化；\chapterref{chap:dedekind-construction}和\chapterref{chap:cauchy-construction}将给出两种构造。

三项目标之间还需要相容性。新加入的点必须可与有理数比较，运算应与有理数原有运算一致，并且不同构造方式得到的边界不应被重复计数。Dedekind 构造以“相同左部”判定同一点，Cauchy 构造以“差为零列”判定同一点。两种等价关系都服务于同一要求：由所有有理精度无法区分的逼近对象应表示同一个数。

此外，完备性是关于数系内部的封闭性。若用一个更大但尚未说明性质的数系作为外部容器，再宣称所有 Cauchy 列在其中有极限，并没有完成实数构造；必须证明极限对象组成有序域，且对其自身的 Cauchy 列仍然封闭。

\subsection{十进制逼近与实数的有限精度表示}


对 \(x\in\R\)、\(n\in\N\)，定义
\[
 d_n(x)=\frac{\lfloor10^n x\rfloor}{10^n}.
\]

\begin{proposition}
对每个 \(x\in\R\) 与 \(n\in\N\)，
\[
 d_n(x)\le x<d_n(x)+10^{-n},
\qquad d_n(x)\le d_{n+1}(x).
\]
\end{proposition}

\begin{proof}
第一式由整数部分定理除以 \(10^n\) 得到。\(d_n(x)\) 是分母为 \(10^{n+1}\) 且不超过 \(x\) 的候选数，而 \(d_{n+1}(x)\) 是此类候选中的最大者，故第二式成立。
\end{proof}

这给出向下截断误差 \(<10^{-n}\)。四舍五入可把绝对误差改进为不超过 \(\frac12 10^{-n}\)，但在恰好位于两个网格点中点时必须另行约定舍入规则。

对负数，\(d_n(x)\) 始终是从下方逼近的网格点。例如
\[
 d_n(-\sqrt2)=-\frac{\lceil10^n\sqrt2\rceil}{10^n},
\]
它通常比“直接删去绝对值的小数尾再加负号”更小。采用向零截断会破坏统一的 \(d_n(x)\le x\) 关系，因此后续上确界论证固定使用向下截断。

\begin{proposition}
对每个 \(x\in\R\)，截断满足
\[
 0\le x-d_n(x)<10^{-n},\qquad
 0\le d_{n+1}(x)-d_n(x)\le9\cdot10^{-(n+1)}.
\]
\end{proposition}

\begin{proof}
第一式已证。相邻网格细分后，\(d_{n+1}(x)-d_n(x)\) 是
\[
 0,10^{-(n+1)},\ldots,9\cdot10^{-(n+1)}
\]
之一，故得第二式。
\end{proof}

\subsection{无限小数表示、循环小数与表示不唯一性}


给定非负整数 \(m\) 和数字 \(a_k\in\{0,\ldots,9\}\)，令
\[
 s_n=m+\sum_{k=1}^n a_k10^{-k}.
\]
部分和集合有上界 \(m+1\)，故可定义无限小数的值为
\[
 m.a_1a_2\cdots:=\sup\{s_n:n\in\Npos\}.
\]
对任意 \(p>n\)，有限等比和给出
\[
 0\le s_p-s_n
 \le \sum_{k=n+1}^{p}9\cdot10^{-k}
 <10^{-n}.
\]
对 \(p\) 取上确界，得到 \(0\le \sup_k s_k-s_n\le10^{-n}\)。这一论证只使用有限和与上确界，没有预先调用无穷级数的定义。

\begin{theorem}[十进制表示]
每个 \(x\ge0\) 都可写成
\[
 x=m+\sup_n\sum_{k=1}^n a_k10^{-k},
 \qquad m\in\N,\quad a_k\in\{0,\ldots,9\}.
\]
若规定数字列不得从某一位起恒为 \(9\)，则该表示唯一。若不作此规定，一个数至多有两种表示；出现两种表示时，其中一种从某一位起全为 \(0\)，另一种在该位向前借一单位并从此全为 \(9\)。借位可以穿过小数点，例如 \(1.000\ldots=0.999\ldots\)。负数的十进制表示按 \(x=-|x|\) 处理，不把负号吸收到整数部分 \(m\) 中。
\end{theorem}

\begin{proof}
取 \(m=\lfloor x\rfloor\)。由
\[
  d_n(x)=\frac{\lfloor10^nx\rfloor}{10^n}
\]
定义整数
\[
  a_n=10^n d_n(x)-10^{n-1}d_{n-1}(x)\qquad(n\ge1).
\]
相邻截断估计说明 \(a_n\in\{0,\ldots,9\}\)，并且
\[
  d_n(x)=m+\sum_{k=1}^n a_k10^{-k}.
\]
由于 \(0\le x-d_n(x)<10^{-n}\)，这些有限小数的上确界为 \(x\)，存在性得证。

设两种表示在整数部分及前 \(k-1\) 位相同，而第 \(k\) 位首次不同，较大一侧该位至少多 \(10^{-k}\)。较小一侧在第 \(k\) 位以后能够补偿的最大量为
\[
  \sup_{p>k}\sum_{j=k+1}^{p}9\cdot10^{-j}=10^{-k}.
\]
所以两者相等只能发生在第 \(k\) 位恰差 \(1\)、较大一侧以后全为 \(0\)、较小一侧以后全为 \(9\) 时。若首次差异发生在整数部分，同一论证就是穿过小数点的借位。禁止最终全为 \(9\) 后，只剩一种表示。
\end{proof}

\begin{theorem}[有理数与循环小数]
取不最终全为 \(9\) 的规范表示。实数为有理数，当且仅当其绝对值的十进制数字从某位起周期重复；终止小数视为以后恒为 \(0\) 的周期表示。
\end{theorem}

\begin{proof}
只需处理 \(x\ge0\)。若 \(x=p/q\)，对小数部分作长除法。每一步余数属于有限集合 \(\{0,1,\ldots,q-1\}\)：余数变为零时表示终止；否则两个余数必相同，而后续除法完全由当前余数决定，故数字开始循环。

反之，设从第 \(N+1\) 位起以长度 \(r\) 循环。令 \(y=10^Nx\)，把其整数部分记为 \(M\)，循环小数部分记为 \(z=y-M\)。数 \(10^rz-z\) 的小数尾逐位抵消，等于由一个循环节组成的整数 \(c\)，所以
\[
  z=\frac{c}{10^r-1},\qquad
  x=10^{-N}\left(M+\frac{c}{10^r-1}\right)\in\Q.
\]
全零循环同样包含在此式中。
\end{proof}

\subsection{实数集不可数}


\begin{theorem}
\(\R\) 不可数。
\end{theorem}

\begin{proof}
\chapterref{chap:mathematical-language}已经证明 \(\mathcal P(\Npos)\) 不可数。对 \(A\subseteq\Npos\) 定义
\[
 \Phi(A)=\sup_n\sum_{k=1}^n\frac{2\chi_A(k)}{3^k}\in[0,1].
\]
若 \(A\ne B\)，令 \(m\) 为二者首次不同的指标，并设 \(m\in A\setminus B\)。则第 \(m\) 位贡献差为 \(2\cdot3^{-m}\)。对每个 \(p>m\)，以后各位的有限反向补偿不超过
\[
 \sum_{k=m+1}^{p}2\cdot3^{-k}<3^{-m}.
\]
对相应有限和取上确界，仍有 \(\Phi(A)-\Phi(B)\ge3^{-m}>0\)。所以 \(\Phi\) 是 \(\mathcal P(\Npos)\) 到 \(\R\) 的单射，\(\R\) 不可数。
\end{proof}

使用三进制数字 \(0,2\) 避开了十进制尾随 \(9\) 的表示歧义。证明还表明不可数性已经出现在区间 \([0,1]\) 内。

另一个常用证明是 Cantor 对角线法：若把 \([0,1]\) 中实数列成 \(x_1,x_2,\ldots\)，统一选取不以 \(9\) 为尾的十进制表示，再构造第 \(n\) 位既不等于 \(x_n\) 的第 \(n\) 位、又只取数字 \(1,2\) 的小数。所得实数与每个 \(x_n\) 至少在一位不同，故不在枚举中。限制新数字避开 \(0,9\) 可以排除双重表示造成的歧义。

\begin{corollary}
每个非退化区间都不可数，而每个区间中的有理数仍至多可数。因此任意非空开区间都含不可数多个无理数。
\end{corollary}

\begin{proof}
仿射双射 \(t\mapsto a+(b-a)t\) 把 \([0,1]\) 映到 \([a,b]\)。任意非退化区间包含某个非退化闭子区间，故不可数。若其中无理数至多可数，则它与至多可数的有理数并仍至多可数，与区间不可数矛盾。
\end{proof}

\section*{习题}
\addcontentsline{toc}{section}{习题}

\subsection*{A组　定义与基本方法}
\begin{enumerate}[label=\textbf{\thechapter.\arabic*.}]
 \exerciseitem{1.1} 证明任意两个不同有理数之间有无穷多个有理数。
 \exerciseitem{1.2} 证明：若整数 \(n^2\) 被素数 \(p\) 整除，则 \(n\) 被 \(p\) 整除；据此证明非平方正整数 \(m\) 的平方根不是有理数。
 \exerciseitem{1.3} 验证正文中关于 \(T(q)\) 的两个恒等式，并分别说明 \(q^2<2\) 与 \(q^2>2\) 时像点所在的位置。
 \exerciseitem{1.4} 证明 \(L\) 向下封闭、\(U\) 向上封闭，并证明二者都没有端点。
 \exerciseitem{1.5} 判断真伪并证明：
 \begin{enumerate}[label=(\arabic*)]
  \item \(S\) 的每个元素都小于 \(S\) 的某个元素；
  \item \(S\) 的上界全体有最小元；
  \item \(S\) 的任意两个上界中较小者仍是上界。
 \end{enumerate}
 \exerciseitem{2.1} 从域公理证明 \(0x=0\)、\((-1)x=-x\) 与 \((-x)(-y)=xy\)。
 \exerciseitem{2.2} 证明 \(x<y\Longleftrightarrow y-x>0\)，并推导不等式加法法则与负数乘法反向法则。
 \exerciseitem{2.3} 证明 \(x^2\ge0\)，且 \(x^2=0\Longleftrightarrow x=0\)。
 \exerciseitem{2.4} 证明对 \(r>0\)，
 \[
 |x-a|<r\Longleftrightarrow a-r<x<a+r,
 \quad
 |x-a|\le r\Longleftrightarrow a-r\le x\le a+r.
 \]
 \exerciseitem{2.5} 证明反三角不等式，并给出取等的简单情形。
 \exerciseitem{2.6} 证明不存在使 \(\C\) 成为有序域的全序。
 \exerciseitem{3.1} 证明对任意 \(x\in\R\)，存在 \(n\in\N\) 使 \(n>|x|\)。
 \exerciseitem{3.2} 证明对任意 \(\varepsilon>0\) 与 \(a>0\)，存在 \(n\in\Npos\) 使 \(a/n<\varepsilon\)。
 \exerciseitem{3.3} 证明整数部分的恒等式
 \[
 \lfloor x+m\rfloor=\lfloor x\rfloor+m\quad(m\in\Z),
 \qquad \lfloor-x\rfloor=-\lceil x\rceil.
 \]
 \exerciseitem{3.4} 证明任意非空开区间含无穷多个有理数和无穷多个无理数。
 \exerciseitem{3.5} 对任意 \(x\in\R\) 与 \(n\in\Npos\)，构造 \(q\in\Q\) 使 \(|x-q|<1/n\)。
\end{enumerate}

\subsection*{B组　证明与反例}
\begin{enumerate}[label=\textbf{\thechapter.\arabic*.},resume]
 \exerciseitem{1.6} 若 \(r\in\Q\)、\(r>0\) 且 \(r^2<2\)，证明存在 \(n\in\Npos\) 使 \((r+1/n)^2<2\)。不得引用实数的 Archimedes 性质。
 \exerciseitem{1.7} 设 \(S_3=\{q\in\Q:q\ge0,q^2<3\}\)。仿照正文证明 \(S_3\) 在 \(\Q\) 中有上界但无上确界。
 \exerciseitem{1.8} 对正文数列 \(a_n\) 证明 \(a_n\le a_{n+1}<a_n+10^{-n}\)，并求 \(a_1,a_2,a_3\)。
 \exerciseitem{1.9} 证明任意在 \(\Q\) 中收敛的有理数列都是 Cauchy 列，并指出证明没有使用完备性。
 \exerciseitem{1.10} 构造一个在 \(\Q\) 中不收敛的有理 Cauchy 列，使其偶数项与奇数项都严格单调。
 \optionalexerciseitem{1.11} 证明两个没有有理极限的有理 Cauchy 列之和可能有有理极限，也可能没有有理极限，各给一例。
 \exerciseitem{2.7} 分别求下列实数集合的上界全体、下界全体、上下确界及最大最小元：
 \begin{enumerate}[label=(\arabic*)]
  \item \((0,1)\)；
  \item \([0,1)\)；
  \item \(\{x\in\R:x\ge0,\ x^2<2\}\)。
 \end{enumerate}
 \exerciseitem{2.8} 证明上确界若存在则唯一，并写出下确界的近似刻画。
 \exerciseitem{2.9} 若 \(\varnothing\ne A\subseteq B\subseteq\R\) 均有上界，证明 \(\sup A\le\sup B\)，并给出严格不等与等号的例子。
 \exerciseitem{2.10} 对非空有上界集合 \(A,B\subseteq\R\)，证明
 \[
 \sup(A+B)=\sup A+\sup B,
 \quad A+B=\{a+b:a\in A,b\in B\}.
 \]
 \exerciseitem{2.11} 对非空有界集合 \(A,B\subseteq\R\)，证明
 \[
 \sup(A\cup B)=\max\{\sup A,\sup B\},\quad
 \inf(A\cup B)=\min\{\inf A,\inf B\}.
 \]
 \optionalexerciseitem{2.12} 构造 \(A,B\subseteq\R\)，使 \(A\cap B\ne\varnothing\) 且
 \[
 \sup(A\cap B)<\min\{\sup A,\sup B\}.
 \]
 \exerciseitem{3.6} 计算 \(d_n(-\sqrt2)\)，并说明向下截断对负数为何不是简单删除小数尾。
 \exerciseitem{3.7} 证明 \(d_n(x)\le d_{n+1}(x)\le d_n(x)+9\cdot10^{-(n+1)}\)。
 \exerciseitem{3.8} 用上确界计算 \(0.999\ldots\) 与 \(0.272727\ldots\)。
 \exerciseitem{3.9} 证明任意最终循环小数为有理数，并给出循环节长度为 \(p\) 时的通式。
\exerciseitem{3.10} 分别证明有限小数全体与最终循环小数全体均为可数无限集。
 \optionalexerciseitem{3.11} 证明代数数的补集在每个非空开区间中不可数。
\end{enumerate}

\subsection*{C组　综合问题}
\begin{enumerate}[label=\textbf{\thechapter.\arabic*.},resume]
 \projectexerciseitem{1.12} \ 【前置：\exerciseref{1.2}；预计用时：45分钟】对非平方正整数 \(m\)，定义 \(L_m=\{q\in\Q:q<0\text{ 或 }q^2<m\}\)。证明它给出无有理端点的分割，并构造分母为 \(10^n\) 的 Cauchy 逼近列。
 \optionalexerciseitem{1.13} 设 \(F\) 为有序域。说明“任意两点之间有第三点”可由域公理推出，因而不能作为完备性公理；再给出一个 \(\Q\) 失败而 \(\R\) 将满足的命题。
 \openexerciseitem{1.14} 比较“补入所有代数数”与“补全所有 Cauchy 列”两种扩张目标。可把\chapterref{chap:real-number-system}的不可数性结论作为条件性说明。
 \exerciseitem{1.15} 证明 \(S\) 的上界全体恰为 \(\{u\in\Q:u>0,u^2>2\}\)，并由此重新推出 \(S\) 无上确界。
 \exerciseitem{1.16} 以任意整数 \(b\ge2\) 为底构造 \(\sqrt2\) 的有理网格逼近列，并证明它是 Cauchy 列。
 \exerciseitem{1.17} 设 \((a_n)\)、\((b_n)\) 分别是底数 \(2\) 与底数 \(10\) 的上述逼近列。证明 \(a_n-b_n\to0\)。
  \exerciseitem{1.18} 证明有理 Cauchy 列有界。再分别给出有界但非 Cauchy 的数列，以及相邻项之差趋于零但非 Cauchy 的数列。
 \optionalexerciseitem{1.19} 设有理 Cauchy 列 \((q_n)\) 满足 \(q_n^2\to2\) 且最终为正。证明任意两条此类列的差趋于零。
 \projectexerciseitem{1.20} \ 【前置：本章全部内容；预计用时：75分钟】围绕集合 \(S\) 建立“分割无端点—无有理上确界—存在无有理极限的 Cauchy 列”的三向联系图；每个箭头给出严格证明并标明是否使用十进制表示。
 \exerciseitem{2.13} 由确界公理推出下确界原理，并证明反向蕴含。
 \exerciseitem{2.14} 在平方根定理中给出 \(h\) 的显式公式，避免使用尚未证明的 Archimedes 性质。
 \exerciseitem{2.15} 证明若 \(a<b\) 与 \(0<t<1\)，则
 \[
 a<(1-t)a+tb<b,
 \]
 并说明该结论不使用完备性。
 \projectexerciseitem{2.16} \ 【前置：\exerciseref{2.8}、2.13；预计用时：60分钟】证明嵌套闭区间 \([a_n,b_n]\) 若满足
 \[
 a_n\le a_{n+1}\le b_{n+1}\le b_n,
 \]
 则交集非空。此处不要求交集为单点。
 \optionalexerciseitem{2.17} 若 \(A\subseteq\R\) 非空有上界且 \(c<\sup A\)，证明存在 \(a\in A\) 使 \(c<a\le\sup A\)。
 \openexerciseitem{2.18} 列出平方根存在性证明中分别使用的域、序、确界三类性质，并说明在 \(\Q\) 中哪一步失败。
 \exerciseitem{2.19} 证明上界集向上封闭、下界集向下封闭，并把 \(\sup A\) 表述为 \(\operatorname{UB}(A)\) 的最小元。
 \exerciseitem{2.20} 证明上确界近似刻画中可把“对每个 \(\varepsilon>0\)”等价改写为“对每个 \(t<s\)”。
 \exerciseitem{2.21} 设 \(A,B\) 非空有上界，且每个 \(a\in A\) 都不超过某个 \(b\in B\)。证明 \(\sup A\le\sup B\)，并构造严格不等与等号例。
 \exerciseitem{2.22} 求 \(\sup\{x/(1+x):x>0\}\)，不得使用微分。
 \exerciseitem{2.23} 设 \(A\) 非空有界，证明
 \[
  \sup\{\abs a:a\in A\}=\max\{\sup A,-\inf A\}.
 \]
 \exerciseitem{2.24} 讨论从非空集合中删除有限个元素后上确界何时保持不变；给出一个充分条件和两个边界反例。
 \optionalexerciseitem{2.25} 设 \(\lambda<0\)。推导 \(\sup(\lambda A)\) 与 \(\inf A\) 的关系，并证明。
 \projectexerciseitem{2.26} \ 【前置：本章确界运算；预计用时：75分钟】建立非空有界集合在平移、正负伸缩、并集、Minkowski 和与乘积下的确界计算表；每个公式注明必要的符号条件。
 \projectexerciseitem{3.12} \ 【前置：\exerciseref{3.7}；预计用时：60分钟】对整数基数 \(b\ge2\) 建立 \(b\) 进制表示，证明误差界、表示不唯一的唯一形式以及“有理数当且仅当最终循环”。
 \exerciseitem{3.13} 在实数不可数性的证明中补全尾和上确界的计算，并证明 \(\Phi\) 的像由只含三进制数字 \(0,2\) 的数构成。
 \optionalexerciseitem{3.14} 证明任意非退化闭区间 \([a,b]\) 与 \([0,1]\) 之间存在双射，并推出每个非退化区间不可数。
 \openexerciseitem{3.15} 给出一个非 Archimedes 有序域的形式模型或参考构造，并指出本章哪些结论首先失效。
 \exerciseitem{3.16} 证明 Archimedes 性质的三种表述：自然数无上界、存在任意小的 \(1/n\)、任意正量的整数倍可超过另一正量，彼此等价。
 \exerciseitem{3.17} 证明
 \[
  \lfloor x\rfloor+\lfloor y\rfloor
  \le\lfloor x+y\rfloor
  \le\lfloor x\rfloor+\lfloor y\rfloor+1,
 \]
 并刻画何时出现右端的 \(+1\)。
 \exerciseitem{3.18} 对任意固定 \(n\) 证明存在 \(m\in\Z\) 使 \(m/n\le x<(m+1)/n\)，并解释这比单纯稠密性多给了什么信息。
 \exerciseitem{3.19} 用 Cantor 对角线法证明 \([0,1]\) 不可数，明确处理十进制双重表示。
 \exerciseitem{3.20} 证明任意非空开区间含不可数多个无理数。
 \exerciseitem{3.21} 设 \(D=\{d_n(x):n\in\N\}\)。求 \(\sup D\)，并讨论 \(x\) 为有限小数时 \(D\) 是否有最大元。
\optionalexerciseitem{3.22} 证明全体实代数数可数无限，并推出超越数在每个非空开区间中不可数。
 \projectexerciseitem{3.23} \ 【前置：本章表示与可数性；预计用时：90分钟】比较二进制、十进制、三进制表示在唯一性、循环性与不可数性证明中的作用；给出一套避免表示歧义的统一约定。
 \openexerciseitem{3.24} 分析浮点数有限集合与实数连续统之间的差别，区分舍入误差、表示误差和运算误差；只讨论数学模型，不依赖具体硬件标准。
\end{enumerate}

\section*{部分习题提示}
\begin{itemize}
 \item \exerciseref{1.6}：直接使 \(1/n<(2-r^2)/(2r+1)\)。
 \item \exerciseref{1.10}：从十进制逼近列抽取严格递增子列。
 \item \exerciseref{1.11}：比较 \(a_n+(-a_n)\) 与 \(a_n+a_n\)。
 \item \exerciseref{1.12}：可用足够小的显式有理增量代替正文的变换 \(T\)。
 \item \exerciseref{2.6}：有序域中平方非负，而 \(i^2=-1\)。
 \item \exerciseref{2.10}：先证右端是上界，再用两个近似元从下方逼近。
\item \exerciseref{2.16}：令 \(s=\sup\{a_n:n\in\Npos\}\)，证明每个 \(b_n\) 都是该集合的上界。
 \item \exerciseref{3.6}：\(\lfloor-10^n\sqrt2\rfloor=-\lceil10^n\sqrt2\rceil\)。
 \item \exerciseref{3.11}：从区间中删去可数代数数，剩余集合不能至多可数。
 \item \exerciseref{3.14}：使用仿射映射 \(t\mapsto a+(b-a)t\)。
\end{itemize}
